% Generated by Claude Opus 5 (Anthropic) — Fatih's paragraph structure + fixes.

\section{Introduction}
\label{sec:intro}

Autonomous cyber-physical systems such as cars, drones and industrial robots run
their software as a chain of stages, typically sensing, perception, planning and
control. The path from a sensor reading to the resulting actuation must complete
within an end-to-end deadline for the system to be considered functional, safe,
or both. Some of these stages require high-throughput computation and routinely
offload it to GPUs. Such systems are usually built on a middleware, which raises
the level of abstraction so that the hundreds of components making up these
stages can be implemented with reasonable effort. The trade-off is visibility:
components are written against the middleware, not against the platform, so the
source code says what the system computes, not when or where it runs.

ROS~2 is one such middleware, widely used in industry and academia. It abstracts
an application into three levels: the callback is the unit of computation; a
node owns a component's interface, the timers, subscriptions, services, clients
and waitables that carry its callbacks together with the publishers it sends
through; and an executor dispatches the ready callbacks of the nodes assigned to
it on one or more threads of an OS process. Three parties decide when and where
all of this runs, and none of them sees the other two: the OS schedules the
executor threads, the executor picks which ready callback runs next, and the GPU
run-time orders the work that callback submits. None of them knows the chain the
callback belongs to---and the chain is what the deadline applies to.

Timing analysis of ROS~2 executors is a recent but fast-establishing field:
there are analyses for the single-threaded
executor~\cite{casini-19-reservationbased}, for the multi-threaded
executor~\cite{arafat-2026-mtexecutor-cbgroups} and for reaction time and data
age~\cite{teper-2024-end2endmtexecutoranalysisoptimization}.
%, and a recent survey reviews the field~\cite{casini-2025-surveyros2}.
A shortcoming common to these analyses is that the model is assumed: the chains,
the callback-to-executor mapping and the callback-group memberships are given as
input, so the analysis describes the timing behaviour of the assumed model, not
of any application that actually runs.

A second line of work recovers the model from a running system.
ROS-Llama~\cite{blass-2021-rosllama-paper} instruments \texttt{rclcpp} and
\texttt{rcl} to recover the callback graph and per-callback execution-time
curves; Abaza et al.~\cite{abaza-2024-ros2modelsynthesis} attach sixteen uprobes
to \texttt{rcl} and \texttt{rmw} entry and exit points and combine them with
\texttt{sched\_switch} to synthesise a callback DAG with measured execution
times; LiMEIT~\cite{brandenburg-2026-limeit} extracts periodic models of ROS~2
callbacks from OS scheduler events and sampled stacks, without framework
knowledge. All three consider only a single-threaded executor. A further family
of tools builds on the instrumentation provided by
\texttt{ros2\_tracing}~\cite{bedard-2022-ros2tracing}.
CARET~\cite{kuboichi-2022-caret} generates an architecture file from the trace
and tracks individual messages, and an extension of it records executors,
callback groups and per-callback timelines~\cite{peng-2022-caretcomp}. All of
them stop at the ROS layer: neither the thread a callback runs on nor the
accelerator it uses is recorded. Neither group covers it, as a recent survey
states: \textit{``the vast majority of the discussed works focus on CPU
scheduling only and neglect the presence of hardware
accelerators''}~\cite{casini-2025-surveyros2}; making accelerators first-class
in a ROS~2 timing model is an open
problem~\cite{yano-2026-rightanalyticalmodelsoftware}.
% Where the accelerator is covered in the ROS~2 context, it is covered as an
% input: PAAM~\cite{enright-2024-paam} and ROSGM~\cite{li-2023-rosgm} route
% accelerator calls through a scheduling server that every accelerator-using
% callback must register with, declaring the accelerator it uses; PAAM's
% admission control additionally takes the worst-case duration of each
% accelerator segment as given.
